% Computer Science A --- SDLC, Agile/Scrum, SOLID/Java bridges, Linux, Maven (\input from \texttt{main.tex}).

\runningheader{Computer Science A}{SDLC, Agile, Linux, \& Maven}{Page \thepage\ of \numpages}

\begingradingrange{csasdlc}

\uplevel{\par\textbf{Software development lifecycle}\par}

\question[4]
Specify the order of the steps in the software development lifecycle. In the table below, write the correct \textbf{phase name} for each position from first (row~1) to last (row~7). (Spelling should match the usual textbook labels.)

\begin{center}
\footnotesize
\renewcommand{\arraystretch}{1.45}
\begin{tabular}{|c|p{0.58\linewidth}|}
  \hline
  \textbf{Order} & \textbf{Lifecycle phase (fill in)} \\
  \hline
  1 & \fillin[Planning][0.52\linewidth] \\
  \hline
  2 & \fillin[Analysis][0.52\linewidth] \\
  \hline
  3 & \fillin[Design][0.52\linewidth] \\
  \hline
  4 & \fillin[Development][0.52\linewidth] \\
  \hline
  5 & \fillin[Testing][0.52\linewidth] \\
  \hline
  6 & \fillin[Deployment][0.52\linewidth] \\
  \hline
  7 & \fillin[Maintenance][0.52\linewidth] \\
  \hline
\end{tabular}
\renewcommand{\arraystretch}{1}
\end{center}
\begin{solution}
  The usual linear ordering is:
  \begin{enumerate}
    \item Planning
    \item Analysis
    \item Design
    \item Development
    \item Testing
    \item Deployment
    \item Maintenance
  \end{enumerate}
  Agile teams still revisit these activities in iterations; the labels describe the classic lifecycle story.
\end{solution}

\question[4]
Which stage of the SDLC usually covers decisions such as choosing technologies and creating UML diagrams?
\begin{choices}
  \choice Build
  \CorrectChoice Design
  \choice Maintain
  \choice Deploy
\end{choices}
\begin{solution}
  \textbf{Design} specifies architecture, interfaces, and major technology choices; UML diagrams (class, sequence, etc.) are typical design artifacts. \textbf{Build/development} implements the design; \textbf{deploy} releases; \textbf{maintain} sustains after release.
\end{solution}

\question[4]
Which of the following best matches a common textbook ordering of \textbf{Software Development Life Cycle} activities?
\begin{choices}
  \CorrectChoice Planning, Requirements, Design, Build, Document, Test, Deploy, Maintain
  \choice Validate, Compile, Test, Package, Integration, Verify, Install, Deploy
  \choice Planning, Code, Test, Design, Build, Maintain, Ship
  \choice Requirements, Planning, Develop, Validate, Deploy, Document, Maintain
\end{choices}
\begin{solution}
  The second list mirrors a \textbf{Maven} lifecycle, not SDLC. The third scrambles order and labels. The fourth permutes planning/requirements. The first option reflects a common expanded linear story (names vary by text: ``Analysis'' vs ``Requirements,'' ``Implementation'' vs ``Build,'' etc.).
\end{solution}

\question[4]
Which of the following is the core difference between the Agile and Waterfall methodologies of software development?
\begin{choices}
  \CorrectChoice Agile is an iterative style and adapts well to change; Waterfall follows a rigid planned schedule
  \choice Waterfall is an iterative style and adapts well to change; Agile follows a rigid planned schedule
  \choice Agile is a team-based approach; Waterfall is used for individual or solo projects
  \choice Agile is used for individual or solo projects; Waterfall is a team-based approach
\end{choices}
\begin{solution}
  \textbf{Waterfall} emphasizes phased, sequential delivery with upfront planning; \textbf{Agile} favors short feedback loops, working increments, and responding to change.
\end{solution}

\newpage

\uplevel{\par\textbf{Agile and Scrum}\par}

\question[4]
Which of these is \emph{not} one of the core role titles on a typical \textbf{Scrum} team?
\begin{choices}
  \CorrectChoice Marketer
  \choice Product Owner
  \choice Developer
  \choice Scrum Master
\end{choices}
\begin{solution}
  Scrum names three accountabilities: \textbf{Product Owner} (value/ordering), \textbf{Developers} (who create the increment), and \textbf{Scrum Master} (coaching the Scrum process). Marketing stakeholders may collaborate, but ``Marketer'' is not one of those three core roles in the framework.
\end{solution}

\question[4]
An Agile workflow \emph{always} produces a better end product than a Waterfall workflow.
\begin{choices}
  \choice TRUE
  \CorrectChoice FALSE
\end{choices}
\begin{solution}
  Outcomes depend on context, skill, requirements stability, and execution. Agile reduces some risks (late feedback, change cost) but does not \textbf{guarantee} a ``better'' product in every situation.
\end{solution}

\question[4]
Why was the \textbf{Agile} model created?
\begin{choices}
  \choice Managers wanted a way to better track employee performance
  \CorrectChoice Allowed clients to see where their end product was at all times
  \choice Developers wanted to put client needs first
  \choice Agile allowed everyone to easily focus on multiple tasks at a time
\end{choices}
\begin{solution}
  Course treatments tie Agile's rise to heavy, slow feedback cycles in plan-driven work: short iterations, demos, and reviews make \textbf{visible progress} and cheaper changes. The Manifesto also stresses customer collaboration and working software; ``multiple tasks at once'' is a distractor (flow/WIP limits are closer to Lean/Kanban). ``Developers put clients first'' is partially true in spirit but less specific than incremental visibility as the quiz's intended contrast to the other distractors.
\end{solution}

\question[4]
Which of the following Agile / Scrum events is \emph{typically} time-boxed to about \textbf{15 minutes}?
\begin{choices}
  \choice Sprint Planning
  \CorrectChoice Daily Stand-up (Daily Scrum)
  \choice Sprint Review
  \choice All of the above (``it's Agile, so everything is 15 minutes'')
\end{choices}
\begin{solution}
  The \textbf{Daily Scrum} is commonly capped at fifteen minutes to synchronize and surface impediments. Sprint Planning, Sprint Review, and Sprint Retrospective are longer, structured events.
\end{solution}

\question[4]
Which artifact is most often used to \textbf{track progress} toward completing the work selected for the current sprint (for example, remaining work over time)?
\begin{choices}
  \choice Sprint Backlog
  \choice UML diagram
  \CorrectChoice Burndown chart
  \choice Product Backlog
\end{choices}
\begin{solution}
  The \textbf{Sprint Backlog} is the team's plan for the sprint (selected Product Backlog items plus an actionable plan). A \textbf{sprint burndown} (or burnup) chart plots remaining work versus time and is the usual teaching answer for ``tracking sprint progress.'' The \textbf{Product Backlog} orders work for the whole product, not sprint burn-down tracking.
\end{solution}

\question[4]
Which of the following is \emph{not} an advantage of the Agile development methodology?
\begin{choices}
  \choice It provides a responsible approach to the development of software since it involves user and customer input
  \choice The customer is involved throughout the development of the product, which leads to better relations and satisfaction
  \CorrectChoice Agile is intended for everyone to focus on one task at a time before moving to the next phase
  \choice When compared to Waterfall, the Agile model has a reduced scope which results in better time allocation and estimation
\end{choices}
\begin{solution}
  The first two are common claimed benefits (collaboration and feedback). The third mis-describes Agile as strict phase-gating with one task at a time---that is closer to a linear/waterfall mental model, not iterative increments. The fourth is at least debatable (Agile does not inherently ``reduce scope''), but the third is the clearest non-advantage because it is a \textbf{false characterization}.
\end{solution}

\newpage

\uplevel{\par\textbf{SOLID and Java types (bridging topics)}\par}

\question[4]
The Open/Closed principle states that classes should be open for modification but closed for extension.
\begin{choices}
  \choice TRUE
  \CorrectChoice FALSE
\end{choices}
\begin{solution}
  \textbf{Open/Closed} (from SOLID) is usually stated the other way: open for \textbf{extension} (new behavior without editing old code) and closed for \textbf{modification} (avoid fragile edits to stable types), often via polymorphism and abstractions.
\end{solution}

\question[4]
What is a valid difference between an \textbf{abstract class} and an \textbf{interface} in Java (typical CS-A simplification)?
\begin{choices}
  \choice Abstract classes are fully abstract while interfaces only have concrete methods
  \choice Abstract classes are fully abstract while interfaces can have abstract and concrete methods
  \CorrectChoice Interfaces are fully abstract while abstract classes can have abstract and concrete methods
  \choice Interfaces are fully abstract while abstract classes can only have concrete methods
\end{choices}
\begin{solution}
  In introductory Java, interfaces specify abstract instance methods (since Java~8, \texttt{default}/\texttt{static} methods complicate the slogan, but exams often keep: interface contracts vs.\ abstract classes that mix abstract and concrete members). Abstract classes need not be ``fully abstract'' (they may provide shared code); they cannot be instantiated when still abstract.
\end{solution}

\newpage

\uplevel{\par\textbf{Linux command line}\par}

\question[4]
The command to create a new directory in Linux is:
\begin{choices}
  \CorrectChoice \texttt{mkdir}
  \choice \texttt{touch}
  \choice \texttt{nano}
  \choice \texttt{vi}
  \choice \texttt{rm}
\end{choices}
\begin{solution}
  \texttt{mkdir} creates directories; \texttt{touch} updates/creates empty files; \texttt{nano}/\texttt{vi} are editors; \texttt{rm} removes files.
\end{solution}

\question[4]
To copy a file in Linux, you use the command:
\begin{choices}
  \CorrectChoice \texttt{cp}
  \choice \texttt{mv}
  \choice \texttt{rm}
  \choice \texttt{cat}
  \choice \texttt{mkdir}
\end{choices}
\begin{solution}
  \texttt{cp} copies; \texttt{mv} moves/renames; \texttt{rm} deletes; \texttt{cat} prints file contents; \texttt{mkdir} makes directories.
\end{solution}

\newpage

\uplevel{\par\textbf{Maven}\par}

\question[4]
Which is the correct format of a minimal POM file?
\begin{choices}
  \CorrectChoice
  \begin{minipage}[t]{0.88\linewidth}
\begin{lstlisting}[style=pomxml]
<project>
    <modelVersion>4.0.0</modelVersion>
    <groupId>com.mycompany.app</groupId>
    <artifactId>my-app</artifactId>
    <version>1</version>
</project>
\end{lstlisting}
  \end{minipage}
  \choice
  \begin{minipage}[t]{0.88\linewidth}
\begin{lstlisting}[style=pomxml]
<project>
    <modelVersion>4.0.0</modelVersion>
    <groupId>com.mycompany.app</groupId>
    <version>1</version>
</project>
\end{lstlisting}
  \end{minipage}
  \choice
  \begin{minipage}[t]{0.88\linewidth}
\begin{lstlisting}[style=pomxml]
<project>
    <modelVersion>4.0.0</modelVersion>
    <artifactId>my-app</artifactId>
    <version>1</version>
</project>
\end{lstlisting}
  \end{minipage}
  \choice None of the above
\end{choices}
\begin{solution}
  A valid project coordinates block includes \texttt{modelVersion}, \texttt{groupId}, \texttt{artifactId}, and \texttt{version}. \texttt{groupId} + \texttt{artifactId} + \texttt{version} identify the artifact (GAV).
\end{solution}

\question[4]
Which is the correct format of dependencies in a \texttt{pom.xml} file?
\begin{choices}
  \choice
  \begin{minipage}[t]{0.88\linewidth}
\begin{lstlisting}[style=pomxml]
<dependency>
    <groupId>.......</groupId>
    <artifactId>.......</artifactId>
</dependency>
\end{lstlisting}
  \end{minipage}
  \CorrectChoice
  \begin{minipage}[t]{0.88\linewidth}
\begin{lstlisting}[style=pomxml]
<dependency>
    <groupId>.......</groupId>
    <artifactId>.......</artifactId>
    <version>......</version>
</dependency>
\end{lstlisting}
  \end{minipage}
  \choice
  \begin{minipage}[t]{0.88\linewidth}
\begin{lstlisting}[style=pomxml]
<dependency>
    <groupId>.......</groupId>
    <artifactId>.......</artifactId>
    <Archetype>.......</Archetype>
    <version>......</version>
</dependency>
\end{lstlisting}
  \end{minipage}
  \choice None of the above
\end{choices}
\begin{solution}
  Each \texttt{dependency} needs \texttt{groupId}, \texttt{artifactId}, and \texttt{version} so Maven can resolve the library. \texttt{Archetype} is not a dependency element (archetypes scaffold new projects).

  The POM (\textbf{Project Object Model}) in \texttt{pom.xml} describes the project and build; coordinates uniquely identify artifacts.
\end{solution}

\question[4]
A project with packaging set to \blank[3em] will produce a typical standalone Java archive artifact (as in ``a Java file'' in the sense of a deployable JAR).
\begin{choices}
  \CorrectChoice JAR
  \choice WAR
  \choice XML
  \choice None of the above
\end{choices}
\begin{solution}
  \textbf{JAR} packages classes and resources for libraries or runnable apps (\texttt{java -jar \ldots}). \textbf{WAR} is a web layout for servlet containers, not the usual standalone ``Java archive'' answer in this stem.
\end{solution}

\question[4]
Which of the following is an advantage of using Maven?
\begin{choices}
  \CorrectChoice Automatically track dependencies, build and deploy your project
  \choice Track source code version history of your project in a central location
  \choice Use Maven annotations to test your project code
  \choice Improve the quality of your code
\end{choices}
\begin{solution}
  Maven resolves dependencies from repositories and runs a standard lifecycle (\texttt{clean}, \texttt{compile}, \texttt{test}, \texttt{package}, \ldots). Version control is Git/SVN, etc.; quality tools may plug in but are not ``what Maven is'' in one line.
\end{solution}

\question[4]
Which of the following is true about the Maven build lifecycle?
\begin{enumerate}
  \item A build lifecycle is a well-defined sequence of phases which define the order in which goals are executed.
  \item A phase represents a stage in the lifecycle.
\end{enumerate}
\begin{choices}
  \choice Only (i) is true
  \choice Only (ii) is true
  \CorrectChoice Both of the above
  \choice None of the above
\end{choices}
\begin{solution}
  Phases (\texttt{validate}, \texttt{compile}, \texttt{test}, \ldots) are ordered stages; plugin goals bind to phases to perform work.
\end{solution}

\question[4]
Which of the following statements about Maven are true?
\begin{choices}
  \choice Only this statement: the development team can automate the project's build infrastructure in almost no time using Maven
  \choice Only this statement: Maven uses a standard directory layout and a default build lifecycle
  \CorrectChoice Both of the above statements are true
  \choice Neither statement is true
\end{choices}
\begin{solution}
  Both claims are standard teaching points: convention over configuration (\texttt{src/main/java}, \texttt{src/test/java}, etc.) and the default lifecycle make new projects quick to build and automate.
\end{solution}

\question[4]
Which phase in the Maven lifecycle runs checks to verify the package is valid and meets quality criteria?
\begin{choices}
  \choice \texttt{integration-test}
  \CorrectChoice \texttt{verify}
  \choice \texttt{install}
  \choice \texttt{deploy}
\end{choices}
\begin{solution}
  \texttt{verify} runs checks (e.g.\ integration tests, quality plugins) after \texttt{package} and before \texttt{install}/\texttt{deploy}.
\end{solution}

\question[4]
Which phase in the Maven lifecycle tests the compiled source code using a suitable unit testing framework?
\begin{choices}
  \choice \texttt{validate}
  \choice \texttt{compile}
  \CorrectChoice \texttt{test}
  \choice \texttt{package}
\end{choices}
\begin{solution}
  The \texttt{test} phase runs unit tests (e.g.\ JUnit) against compiled production and test classes.
\end{solution}

\question[4]
Which phase in the Maven lifecycle copies the final package to the \textbf{remote} repository for sharing with other developers and projects?
\begin{choices}
  \choice \texttt{integration-test}
  \choice \texttt{verify}
  \choice \texttt{install}
  \CorrectChoice \texttt{deploy}
\end{choices}
\begin{solution}
  \texttt{deploy} publishes to a remote repository; \texttt{install} only puts the artifact in the \textbf{local} repository.
\end{solution}

\question[4]
Which phase in the Maven lifecycle installs the package into the \textbf{local} repository for use as a dependency in other projects on your machine?
\begin{choices}
  \choice \texttt{integration-test}
  \choice \texttt{verify}
  \CorrectChoice \texttt{install}
  \choice \texttt{deploy}
\end{choices}
\begin{solution}
  \texttt{install} copies the built artifact into \texttt{\textasciitilde/.m2/repository} (by default) so other local builds can depend on it.
\end{solution}

\newpage

\uplevel{\par\textbf{Java vocabulary}\par}

\question[4]
Which of these does \emph{not} describe Java (in the usual CS-A sense)?
\begin{choices}
  \choice Object-oriented
  \choice Statically typed
  \choice Platform-independent (via the JVM)
  \CorrectChoice Can run in the browser (like typical client-side web scripts)
\end{choices}
\begin{solution}
  Java bytecode runs on the JVM; it is not the default language for browser scripting (that role belongs to JavaScript). Historical Java applets are obsolete for practical purposes.
\end{solution}

\question[4]
A parameter is a \blank[3.5em] that is passed into a \blank[3.5em] and declared in the \blank[3.5em].
\begin{choices}
  \choice \texttt{method}; \texttt{class}; \texttt{constructor}
  \choice \texttt{variable}; \texttt{package}; \texttt{class}
  \choice \texttt{annotation}; \texttt{class}; \texttt{package}
  \CorrectChoice \texttt{variable}; \texttt{method}; \texttt{method signature}
\end{choices}
\begin{solution}
  Parameters are variables declared in the method signature; arguments are the values supplied at the call site.
\end{solution}

\endgradingrange{csasdlc}
